\documentclass[11pt]{article}
\usepackage[margin=0.9in]{geometry}
\usepackage{graphicx,amsmath,amssymb,booktabs,hyperref,xcolor}
\graphicspath{{/tmp/dd_k3_method_c_v2_figs/}}
\hypersetup{colorlinks=true,linkcolor=blue,urlcolor=blue}
\setlength{\parskip}{6pt}
\title{Method C v2: Bayesian-symmetric monotone ansatz ladder\\
       (structural deficit across $\tau$)}
\author{Fixed-Point Factory --- Method C v2 ladder}
\date{June 2026}

\begin{document}
\maketitle

\section*{What this measures}

For each $\tau$, fit the parametric family
$$
P(u_1, u_2, u_3) = \sigma\big(h(u_1) + h(u_2) + h(u_3)\big)
$$
where $h$ is odd and increasing ($h(u) = \int_0^u \sigma_*^2(s)\,ds$ with
$\sigma_*$ an even Cheby polynomial of degree 8). 5 unknowns. The fit
residual $|F|_\infty = \max|\Phi[P] - P|$ measures \textbf{how much the
equilibrium needs cross-signal joint structure beyond the additive form}.
This is the structural deficit, separated from kernel-band bias.

\section*{Residual along the $\tau$ ladder}

\begin{center}
\includegraphics[width=\linewidth]{fig1_F.png}
\end{center}

\begin{center}
\includegraphics[width=\linewidth]{fig2_loglog.png}
\end{center}

\section*{P range}

\begin{center}
\includegraphics[width=\linewidth]{fig3_Prange.png}
\end{center}

\section*{$\sigma$ coefficients along ladder}

\begin{center}
\includegraphics[width=\linewidth]{fig4_coefs.png}
\end{center}

\section*{Diagnostics}

\begin{center}
\includegraphics[width=\linewidth]{fig5_diag.png}
\end{center}

\section*{Reproducibility}

\texttt{projects/REZN/solver\_code/highprec\_dd\_qd/dd\_k3\_method\_c\_v2.py}
(operator), \texttt{dd\_k3\_method\_c\_v2\_ladder.py} (ladder),
\texttt{dd\_k3\_method\_c\_v2\_figs.py} (figures). Per-tau coefs:
\texttt{projects/REZN/solved\_fixed\_points/dd\_k3\_overnight/method\_c\_v2/coefs/}.

\end{document}
